\documentclass[UKenglish,usenames,dvipsnames,svgnames,table,aspectratio=169,mathserif]{beamer}

\mode<presentation> {

%\usetheme{default}
\usetheme{Madrid}

\setbeamertemplate{footline} % To remove the footer line in all slides uncomment this line

\setbeamertemplate{navigation symbols}{} % To remove the navigation symbols from the bottom of all slides uncomment this line
}

\usepackage{graphicx} % Allows including images
\usepackage{booktabs} % Allows the use of \toprule, \midrule and \bottomrule in tables
\usepackage{hyperref}
\usepackage{apacite}
\usepackage{babel}
\usepackage{fancyvrb}
\usepackage{color}
\usepackage{alltt}
\usepackage{listings}
\usepackage{framed}
\usepackage{courier}
\usepackage{minted}
\usepackage{epstopdf}
\usepackage{xifthen}
\usepackage[utf8]{inputenc}
\usepackage[T1]{fontenc}
\usepackage{textcomp}
\usepackage{gensymb}
\usepackage{svg}
\usepackage{pdfpages}
\usepackage{isodate}

\hypersetup{colorlinks=false}

\setbeamertemplate{bibliography entry title}{}
\setbeamertemplate{bibliography entry location}{}
\setbeamertemplate{bibliography entry note}{}
\setbeamertemplate{itemize items}[circle]
\setbeamertemplate{enumerate items}[circle]
\beamertemplatenavigationsymbolsempty
\setbeamertemplate{footline}{}

\newminted{haskell}{}

\definecolor{g}{RGB}{0,100,0}
\newcommand{\highlight}[1]{\colorbox{yellow}{#1}}
\newcommand{\nega}[1]{\colorbox{yellow}{#1}}
\newcommand{\posi}[1]{\colorbox{green}{#1}}
\newcommand{\nl}{\vspace{\baselineskip}}
\newcommand{\pnl}{\pause \nl}

\graphicspath{{diagrams/}}

\newcommand{\textslide}[1]{{
\begin{frame}
\begin{center}

#1

\end{center}
\end{frame}
}}

\newcommand{\textslideleft}[1]{{
\begin{frame}

#1

\end{frame}
}}

\newcommand{\codeslide}[1]{{
\begin{frame}[fragile]
\begin{haskellcode}
#1
\end{haskellcode}
\end{frame}
}}

\newcommand{\imageslide}[2][1]{{
\begin{frame}\begin{center}
\includegraphics[scale=#1]{#2}
\end{center}\end{frame}
}}

\newcommand{\imageslideleft}[2][1]{{
\begin{frame}
\includegraphics[scale=#1]{#2}
\end{frame}
}}

\newcommand{\imagetextslide}[3][1]{{
\begin{frame}\begin{center}

{#3}

\includegraphics[scale=#1]{#2}
\end{center}\end{frame}
}}

%%----------------------------------------------------------------------------------------
%	TITLE PAGE
%----------------------------------------------------------------------------------------

\title[Monoids]{Monoids: an introduction to lawful structures}
\author{George Wilson}
\institute[]
{
\href{george@wils.online}{george@wils.online}
}

\selectlanguage{UKenglish}
\date{\printdate{2026-10-13}}

\begin{document}

\begin{frame}
\titlepage
\end{frame}


\begin{frame}
\frametitle{Outline}
\begin{center}
\begin{minipage}{0.6\textwidth}
Let's look at:
\begin{itemize}
\item Basic Haskell syntax
\item A simple example type class
\item The Monoid type class (including laws)
\item Examples of monoid instances
\item Examples of functions defined for all monoids
\item Why laws are helpful
\end{itemize}
\end{minipage}
\end{center}
\end{frame}


\begin{frame}[fragile]
\frametitle{Values and types}
\begin{center}
\begin{minipage}{0.6\textwidth}
\begin{haskellcode}
x :: Integer
x = 5 + 10
\end{haskellcode}
\end{minipage}
\end{center}
\end{frame}


\begin{frame}[fragile]
\frametitle{Strings and concatenation}
\begin{center}
\begin{minipage}{0.6\textwidth}
\begin{haskellcode}
bfpg :: String
bfpg = "BFPG"

greeting :: String
greeting = "Hello, " ++ bfpg ++ "!"

\end{haskellcode}
\pause
\hrule
\begin{haskellcode}

> greeting
-- "Hello, BFPG!"
\end{haskellcode}
\end{minipage}
\end{center}
\end{frame}


\begin{frame}[fragile]
\frametitle{Functions on strings}
\begin{center}
\begin{minipage}{0.6\textwidth}
\begin{haskellcode}
greet :: String -> String
greet name = "Hello, " ++ name ++ "!"

\end{haskellcode}
\pause
\hrule
\begin{haskellcode}

> greet bfpg
-- "Hello, BFPG!"
\end{haskellcode}
\end{minipage}
\end{center}
\end{frame}


\begin{frame}[fragile]
\frametitle{Custom datatypes}
\begin{center}
\begin{minipage}{0.6\textwidth}
\begin{haskellcode}
data MeetupGroup =
  MeetupGroup {
    name :: String,
    url :: String,
    description :: String
  }
\end{haskellcode}
\pause
\begin{haskellcode}
bfpgMeetup :: MeetupGroup
bfpgMeetup =
  MeetupGroup {
    name = "Brisbane Functional Programming Group",
    url = "https://luma.com/bfpg",
    description = "a bunch of lambda heads"
  }
\end{haskellcode}
\end{minipage}
\end{center}
\end{frame}


\begin{frame}[fragile]
\frametitle{Custom data types}
\begin{center}
\begin{minipage}{0.6\textwidth}
\begin{haskellcode}
data TrafficLight
  = Red
  | Yellow
  | Green
\end{haskellcode}
\end{minipage}
\end{center}
\end{frame}


\begin{frame}[fragile]
\frametitle{Type classes}
\begin{center}
\begin{minipage}{0.6\textwidth}
\begin{haskellcode}
class Show a where
  show :: a -> String

\end{haskellcode}
\pause
\begin{haskellcode}
instance Show TrafficLight where
  show tl =
    case tl of
      Red -> "Red"
      Yellow -> "Yellow"
      Green -> "Green"
\end{haskellcode}
\end{minipage}
\end{center}
\end{frame}


\begin{frame}[fragile]
\frametitle{Functions with type class constraints}
\begin{center}
\begin{minipage}{0.6\textwidth}
\begin{haskellcode}
greet :: Show a => a -> String
greet s = "Hello, " ++ show s ++ "!"

\end{haskellcode}
\pause
\hrule
\begin{haskellcode}

> greet Yellow
-- "Hello, Yellow!"
\end{haskellcode}
\pause
\begin{haskellcode}

> greet 5
-- "Hello, 5!"
\end{haskellcode}

\pause
\begin{haskellcode}

> greet (++)
-- Error: No Show instance for
--        (String -> String -> String)
\end{haskellcode}

\end{minipage}
\end{center}
\end{frame}


\textslide{
We've seen Haskell expressions, data types, functions, and type classes.
\medskip

Any questions so far?
}

\begin{frame}[fragile]
\frametitle{Monoid}
\begin{center}
\begin{minipage}{0.6\textwidth}
\begin{haskellcode}
class Monoid m where
  mappend :: m -> m -> m
  mempty  :: m
\end{haskellcode}
\pause
\begin{haskellcode}
instance Monoid String where
  mappend s t = s ++ t
  mempty = ""
\end{haskellcode}
\pause
\begin{haskellcode}
-- mappend as an infix operator, for convenience
x <> y = mappend x y
\end{haskellcode}
\end{minipage}
\end{center}
\end{frame}


\begin{frame}[fragile]
\frametitle{Monoid laws}
\begin{center}
\begin{minipage}{0.6\textwidth}
Left identity:

\begin{haskellcode}
-- for all x.
mempty <> x == x
\end{haskellcode}

\pause
\hrule
\bigskip

Right identity:

\begin{haskellcode}
-- for all x.
x <> mempty == x
\end{haskellcode}

\pause
\hrule
\bigskip

Associativity:

\begin{haskellcode}
-- for all x y z.
x <> (y <> z) == (x <> y) <> z
\end{haskellcode}

\end{minipage}
\end{center}
\end{frame}


\begin{frame}[fragile]
\frametitle{Monoid laws specialised to strings}
\begin{center}
\begin{minipage}{0.6\textwidth}
Left identity:

\begin{haskellcode}
-- for all x.
"" ++ x == x
\end{haskellcode}

\hrule
\bigskip

Right identity:

\begin{haskellcode}
-- for all x.
x ++ "" == x
\end{haskellcode}

\hrule
\bigskip

Associativity:

\begin{haskellcode}
-- for all x y z.
x ++ (y ++ z) == (x ++ y) ++ z
\end{haskellcode}

\end{minipage}
\end{center}
\end{frame}


\begin{frame}[fragile]
\frametitle{Pretty printing}
\begin{center}
\begin{minipage}{0.6\textwidth}
Imagine we are trying to pretty print LISP s-expressions
\begin{haskellcode}




expr =
  openParen <> "+" <> space <> "3"
    <> space <> "4" <> closeParen

-- (+ 3 4)
\end{haskellcode}
\end{minipage}
\end{center}
\end{frame}


\begin{frame}[fragile]
\frametitle{Pretty printing}
\begin{center}
\begin{minipage}{0.6\textwidth}
1: Factor out the parens to a helper function

\begin{haskellcode}
parens x = openParen <> x <> closeParen



expr =
  parens ("+" <> space <> "3" <> space <> "4")


-- (+ 3 4)
\end{haskellcode}
\end{minipage}
\end{center}
\end{frame}


\begin{frame}[fragile]
\frametitle{Pretty printing}
\begin{center}
\begin{minipage}{0.6\textwidth}
2: Factor out the space separation to a helper function

\begin{haskellcode}
parens x = openParen <> x <> closeParen

spaceSep list = ... -- omitted for brevity

expr =
  parens (spaceSep ["+", "3", "4"])


-- (+ 3 4)
\end{haskellcode}
\end{minipage}
\end{center}
\end{frame}


\textslide{
Extracting a helper function usually depends on associativity

\medskip
Extracting helper functions is awesome

\medskip
Therefore associativity is awesome!
}


\textslide{What are some other types that fit monoid?}


\begin{frame}[fragile]
\frametitle{Monoid laws specialised to addition}
\begin{center}
\begin{minipage}{0.6\textwidth}
Left identity:

\begin{haskellcode}
-- for all x.
0 + x == x
\end{haskellcode}

\hrule
\bigskip

Right identity:

\begin{haskellcode}
-- for all x.
x + 0 == x
\end{haskellcode}

\hrule
\bigskip

Associativity:

\begin{haskellcode}
-- for all x y z.
x + (y + z) == (x + y) + z
\end{haskellcode}

\end{minipage}
\end{center}
\end{frame}


\begin{frame}[fragile]
\frametitle{Monoid laws specialised to multiplication}
\begin{center}
\begin{minipage}{0.6\textwidth}
Left identity:

\begin{haskellcode}
-- for all x.
1 * x == x
\end{haskellcode}

\hrule
\bigskip

Right identity:

\begin{haskellcode}
-- for all x.
x * 1 == x
\end{haskellcode}

\hrule
\bigskip

Associativity:

\begin{haskellcode}
-- for all x y z.
x * (y * z) == (x * y) * z
\end{haskellcode}

\end{minipage}
\end{center}
\end{frame}


\begin{frame}[fragile]
\frametitle{Integer instances?}
\begin{center}
\begin{minipage}{0.6\textwidth}

We can't have multiple instances like this in Haskell

\begin{haskellcode}
instance Monoid Integer where
  mappend x y = x + y
  mempty = 0

instance Monoid Integer where
  mappend x y = x * y
  mempty = 1
\end{haskellcode}
\end{minipage}
\end{center}
\end{frame}


\begin{frame}[fragile]
\frametitle{Integer instances with newtype wrappers}
\begin{center}
\begin{minipage}{0.6\textwidth}
\begin{haskellcode}
newtype Sum = Sum { unwrapSum :: Integer }
\end{haskellcode}
\pause

\begin{haskellcode}
instance Monoid Sum where
  mappend (Sum x) (Sum y) = Sum (x + y)
  mempty = Sum 0
\end{haskellcode}
\pause
\bigskip

\begin{haskellcode}
newtype Product = Product { unwrapProduct :: Integer }
instance Monoid Product where
  mappend (Product x) (Product y) = Product (x * y)
  mempty = Product 1
\end{haskellcode}
\end{minipage}
\end{center}
\end{frame}


\textslide{Now we have a bunch of monoids.

What can we do with them?}


\begin{frame}[fragile]
\frametitle{Derived functions}
\begin{center}
\begin{minipage}{0.6\textwidth}
\begin{haskellcode}
mconcat :: Monoid m => [m] -> m

\end{haskellcode}

\pause
\hrule

\begin{haskellcode}
> mconcat [Sum 1, Sum 2, Sum 3, Sum 4]
-- Sum 10
\end{haskellcode}

\pause

\begin{haskellcode}
> mconcat (map Sum [1, 2, 3, 4])
-- Sum 10
\end{haskellcode}

\pause

\begin{haskellcode}
> mconcat (map Sum [])
-- Sum 0
\end{haskellcode}
that's mempty
\end{minipage}
\end{center}
\end{frame}


\begin{frame}[fragile]
\frametitle{Derived functions}
\begin{center}
\begin{minipage}{0.6\textwidth}

Is mconcat implemented using a left fold or a right fold?

\begin{haskellcode}
Sum 1 <> (Sum 2 <> (Sum 3 <> (Sum 4 <> mempty)))
\end{haskellcode}

or

\begin{haskellcode}
(((mempty <> Sum 1) <> Sum 2) <> Sum 3) <> Sum 4
\end{haskellcode}

\bigskip

Answer: We don't have to care, because of associativity

\end{minipage}
\end{center}
\end{frame}


\begin{frame}
\frametitle{Parallelism}
\begin{center}
\includesvg[scale=0.6]{diagrams/multicore1.drawio.svg}
\end{center}
\end{frame}
\begin{frame}
\frametitle{Parallelism}
\begin{center}
\includesvg[scale=0.6]{diagrams/multicore2.drawio.svg}
\end{center}
\end{frame}
\begin{frame}
\frametitle{Parallelism}
\begin{center}
\includesvg[scale=0.6]{diagrams/multicore3.drawio.svg}
\end{center}
\end{frame}



\begin{frame}[fragile]
\frametitle{Putting Monoids together}
\begin{center}
\begin{minipage}{0.6\textwidth}
\begin{haskellcode}
data Mean =
  Mean {
    total :: Integer,
    count :: Integer
  }
\end{haskellcode}
\pause

\begin{haskellcode}
instance Monoid Mean where
  mappend (Mean t1 c1) (Mean t2 c2) =
    Mean (t1 + t2) (c1 + c2)
  mempty =
    Mean 0 0
\end{haskellcode}
\pause

\begin{haskellcode}
makeMean :: Integer -> Mean
makeMean num = Mean { total = num, count = 1 }

getMean :: Mean -> Double
getMean (Mean total count) =
  fromIntegral total / fromIntegral count
\end{haskellcode}
\end{minipage}
\end{center}
\end{frame}


\begin{frame}[fragile]
\frametitle{Means}
\begin{center}
\begin{minipage}{0.6\textwidth}
\begin{haskellcode}
> mconcat (map makeMean [1,2,3,4,5])
-- Mean { total = 15, count = 5 }
\end{haskellcode}
\pause
\begin{haskellcode}
> getMean (mconcat (map makeMean [1,2,3,4,5]))
-- 3.0
\end{haskellcode}
\end{minipage}
\end{center}
\end{frame}


\begin{frame}[fragile]
\frametitle{Generalised mconcat}
\begin{center}
\begin{minipage}{0.6\textwidth}
\begin{haskellcode}
class Foldable f where
  foldMap :: Monoid m => (a -> m) -> f a -> m
\end{haskellcode}
\pause
\begin{haskellcode}
instance Foldable [] where
  foldMap a2m list = mconcat (map a2m list)
\end{haskellcode}
\end{minipage}
\end{center}
\end{frame}


\begin{frame}[fragile]
\frametitle{Generalised mean}
\begin{center}
\begin{minipage}{0.6\textwidth}
\begin{haskellcode}
mean :: Foldable f => f Integer -> Double
mean ints = getMean (foldMap makeMean ints)
\end{haskellcode}
\medskip

This function takes the mean of a structure, without caring about the "shape" of the structure
\end{minipage}
\end{center}
\end{frame}


\begin{frame}[fragile]
\frametitle{Trees}
\begin{center}
\begin{minipage}{0.6\textwidth}
\begin{haskellcode}
data Tree a = Leaf a | Branch (Tree a) (Tree a)

\end{haskellcode}

\pause
\begin{haskellcode}
instance Foldable Tree where
  foldMap a2m tree =
    case tree of
      Leaf a -> a2m a
      Branch t u ->
        foldMap a2m t <> foldMap a2m u
\end{haskellcode}
\end{minipage}
\end{center}
\end{frame}


\begin{frame}[fragile]
\frametitle{Trees}
\begin{columns}
\begin{column}{0.5\textwidth}
\begin{haskellcode}
myTree =
  Branch
    (Branch (Leaf 1) (Leaf 2))
    (Branch
      (Branch (Leaf 3) (Leaf 4))
      (Leaf 5)
    )
\end{haskellcode}
\end{column}
\begin{column}{0.5\textwidth}
\includesvg[scale=0.4]{diagrams/tree1.drawio.svg}
\end{column}
\end{columns}
\end{frame}


\begin{frame}[fragile]
\frametitle{Trees}
\begin{columns}
\begin{column}{0.5\textwidth}
\begin{haskellcode}
foldMap Sum myTree
\end{haskellcode}
\end{column}
\begin{column}{0.5\textwidth}
\includesvg[scale=0.4]{diagrams/tree2.drawio.svg}
\end{column}
\end{columns}
\end{frame}


\textslide{There are many more foldable structures}


\begin{frame}
\frametitle{Outline}
\begin{center}
\begin{minipage}{0.6\textwidth}
We've seen:

\begin{itemize}
\item Monoid typeclass and its laws
\item multiple instances of monoid (String, Sum, Product)
\item Laws allowed us to extract helper functions
\item combining instances to make new instances (Mean)
\item Functions that work on all monoids (mconcat, foldMap)
\item Laws allowed us to concatenate structures regardless of their shapes
\end{itemize}
\end{minipage}
\end{center}
\end{frame}


\textslide{
There is much more to monoids.
\medskip

There are many other sturctures with laws.
}

\textslide{\Large{Thanks for listening!}}


\begin{frame}[fragile]
\frametitle{spaceSep}
\begin{center}
\begin{minipage}{0.6\textwidth}
\begin{haskellcode}
spaceSep [] = mempty
spaceSep (x:[]) = x
spaceSep (x:y:rest) = x <> space <> spaceSep (y:rest)
\end{haskellcode}
\end{minipage}
\end{center}
\end{frame}

\end{document}
